\section{Analytic exclusion of the unit diagonal}\label{app:unit}
The classical rule is: read white/right or black/left, turn, toggle the current
cell, and move one edge on the entire square lattice. The initial black set is
finite. Suppose a nonempty word $u$ repeats forever, restores the heading, and
has nonzero full displacement $d$. Set $P=|u|$ and $g=\#R-\#L$. The period
need not be least; arbitrary unread finite debris is allowed.

\begin{theorem}\label{ud:thm:unit}
No such periodic highway has full displacement $(1,-1)$ or any quarter-turn
rotation of this vector, regardless of period, width, growth, or finite seed.
\end{theorem}
The full vector is essential: standard P104 has displacement $(2,-2)$ and is
not excluded because its primitive direction is $(1,-1)$.

We prove the exclusion by forced layer words and a binary interleaving
contradiction. This supplies an analytic alternative to the retained finite
certificate.

\subsection{Stable ranks and the singleton extreme}
Write the phase positions as $q_i$, $0\le i<P$. For $d=(1,-1)$, exactly one
full translation class lies on each level $t=x+y$. Choose $b_t=(0,t)$ and write
$q_i=b_t+a_i d$. By \eqref{eq:stable-order} and Lemma~\ref{lem:criterion},
increasing stable keys $\tau_i=i-Pa_i$ give ranks $k$ whose turns alternate
$R,L,R,L,\ldots$.
A fixed path of $\ell$ moves and displacement $\delta d$ changes
these stable times by $\ell-P\delta$.

The invariant $\sigma$ is $(-1)^{x+y}$ for vertical incoming headings and its
negative for horizontal incoming headings. Every step changes both factors.
At a transverse extreme, the incoming axis and the inward-neighbor condition
fix both directions of a visit. All turns of that class are identical, so
stable alternation forces one singleton $R$ event. Choose the top event $T$
as a sufficiently late phase origin and translate it to $(0,0)$. It arrives
east, leaves south, and gives $\sigma=-1$. The bottom arrives west and leaves
north and must also have even transverse level. The width $W$ is therefore
positive and even. The new phase's complete field remains finite.

We will repeatedly use two direct consequences of chronology.

\paragraph{Forced matching.}
If a directed edge satisfies $b_s+e=b_t+\delta d$, it changes stable time by
$1-P\delta$. All source occurrences and all target occurrences have the same
order and equal cardinality. They match first to first, not by an arbitrary
permutation. At the period seam, the additional epoch cancels in $i-Pa_i$.

\paragraph{Consecutive-return rule.}
A fixed contiguous path's first return to its starting physical cell must
increase that cell's stable rank by exactly one. Translate the path far enough
along $d$ for all its reads to be stable. Any intervening read would be another
node of the same contiguous path; if none occurs, no stable rank can be skipped.
If the endpoints were the same phase event, the path would instead form a
closed component. Such a component cannot occur while obligatory deeper events
remain: the successor on the $P$ phase indices is one cycle, including when
the word is a power of a shorter word.

\subsection{The first two interior layers are forced}\label{ud:sec:forced}
Build downwards. Let $a_k=1$ indicate arrival from the already constructed
upper side and $b_k=1$ departure to that side. At even $t$ the incoming headings
for bits $1,0$ are W,E, and at odd $t$ they are S,N. Literal turning gives
\begin{equation}\label{ud:eq:port}
 b_k=a_k\xor(k\bmod2)\xor(t\bmod2).
\end{equation}
If the word has length $n$, write $I=\sum a_k$, $J=\sum b_k$, $o=n\bmod2$,
and let $e$ count arrival-ones at even ranks. Counting departures gives
\begin{equation}\label{ud:eq:count}
 n=\begin{cases}2(I+J-2e)+o,&t\text{ even},\\
                 2(J-I+2e)-o,&t\text{ odd}.
    \end{cases}
\end{equation}
The remaining lower cut counts are $n-J$ and $n-I$. They agree, starting
with $I=J=1$ from the top singleton.

Width zero is impossible because a move changes $t$. For $W=2$, the extreme
singletons force two middle visits, whereas~\eqref{ud:eq:count} gives $2=4e-o$;
since $n=2$ implies $o=0$, this is impossible. Thus $W\ge4$, and layers
$-1,-2,-3$ are strictly interior, with a deeper visited layer still required.

\subsubsection{First layer}
Denote rank $k$ at level $-1$ by $A_k$. Here $I=J=1$, so the only feasible
arrival words are
\[
 100,\quad1000,\quad001,\quad0010.
\]
Indeed~\eqref{ud:eq:count} gives $n=4e-o$, and its single arrival-one lies at
even rank zero or two. For either $001$ or $0010$, where the arrival-one has rank two, forced matching gives
$A_0\to T\to A_2$, with directions E,S. This travels by $d$ in two updates,
so $\tau(A_2)-\tau(A_0)=2-P<0$, contrary to increasing rank. There are at
least four distinct phase events, so $P\ge4$. Only $100$ and $1000$ remain.
Both have the fixed top path $A_2\to T\to A_0$.

\subsubsection{Second layer}
Let $B_k$ denote rank $k$ at level $-2$. Let $p_0<p_1<\cdots$ list its
arrival-one ranks and $q_0<q_1<\cdots$ its departure-one ranks. Each list has
size $c=2$ or $3$, according as A has length three or four. Matching forces
\[
 B_{q_1}\ \xrightarrow{\mathrm N}\ A_2
 \ \xrightarrow{\mathrm E}\ T
 \ \xrightarrow{\mathrm S}\ A_0
 \ \xrightarrow{\mathrm W}\ B_{p_0}.
\]
This is a first physical return, so the consecutive-return rule gives
\begin{equation}\label{ud:eq:first}
 p_0=q_1+1.
\end{equation}
At this even layer, every even rank belongs to both lists or neither, and
every odd rank belongs to exactly one. Before $p_0$ there is no arrival-one.
Thus every even rank there is in neither list, and every odd rank is a
departure-one. Equation~\eqref{ud:eq:first} forces
\[
 q_0=1,\qquad q_1=3,\qquad p_0=4.
\]
The even rank four belongs to both lists, giving $q_2=4$. Therefore $c=3$:
A has word $1000$. The departure list is exhausted; all later even ranks
have arrival bit zero and all later odd ranks arrival bit one. The remaining
arrivals must be at five and seven, and rank nine cannot occur. Hence B is
exactly one of
\begin{equation}\label{ud:eq:B}
 00001101,\qquad000011010.
\end{equation}
No case enumeration is used in this forcing argument.

The fixed upper paths now include
\begin{align*}
 &B_1\to A_1\to B_5,\\
 &B_3\to A_2\to T\to A_0\to B_4\to A_3\to B_7.
\end{align*}
The B ranks departing downwards are $0,2,5,6,7$ and optionally eight; those
arriving from below are $0,1,2,3,6$ and optionally eight. Only the first five
of each list will be used, so the optional last event creates no further case.

\subsection{A third-layer interleaving contradiction}\label{ud:sec:interleaving}
Let $C_k$ denote rank $k$ at level $-3$ and reuse $p_j,q_j$ for its arrival-one
and departure-one ranks. There are at least five of each. Ordered matching
forces the following paths:
\begin{align*}
 C_{q_1}&\to B_1\to A_1\to B_5\to C_{p_2},\\
 C_{q_3}&\to B_3\to A_2\to T\to A_0\to B_4\to A_3\to B_7\to C_{p_4}.
\end{align*}
Their compass words are respectively E,N,W,S and E,N,E,S,W,N,W,S. Both are
first returns to the starting physical cell, with no other third-layer read
between. Thus
\begin{equation}\label{ud:eq:returns}
 p_2=q_1+1,\qquad p_4=q_3+1.
\end{equation}
Two bridges, however, are forbidden. If $p_1=q_1$, then
\[
 B_2\to C_{p_1=q_1}\to B_1\to A_1\to B_5
\]
has compass word S,E,N,W: its first physical return skips from rank two to
five on B. If $p_4=q_4$, then
\[
 B_4\to A_3\to B_7\to C_{p_4=q_4}\to B_6
\]
has compass word N,W,S,E: its first return skips from rank four to six.
Both violate consecutivity, so
\begin{equation}\label{ud:eq:bridges}
 p_1\ne q_1,\qquad p_4\ne q_4.
\end{equation}

\begin{lemma}[Binary interleaving]\label{ud:lem:binary}
Let $p$ and $q$ be increasing lists of indices in $\{0,\ldots,n-1\}$.
Suppose every even index in this interval belongs to exactly one list, and
every odd index to both or neither. If $p_2=q_1+1$ and $p_4=q_3+1$,
then $p_1=q_1$ or $p_4=q_4$.
\end{lemma}
\begin{proof}
If $q_1$ is odd, it also lies in $p$, and the next index is $p_2$, so
$q_1=p_1$. Otherwise $q_1$ is even and $p_2$ is odd, giving $p_2=q_2$.
If $q_3$ is even, $p_4$ is odd and gives $p_4=q_4$. Otherwise $q_3$ is odd
and is $p_3$. Avoiding both asserted equalities therefore gives
$p_2=q_2$ and $p_3=q_3$, two distinct odd indices that are consecutive in
both lists. No list member can lie strictly between them. But an even
integer lies between distinct odd integers, and it must belong to exactly
one list, a contradiction.
\end{proof}
At odd layer $-3$, equation~\eqref{ud:eq:port} is precisely the list condition
of Lemma~\ref{ud:lem:binary}. Thus~\eqref{ud:eq:returns} contradicts~\eqref{ud:eq:bridges}.

\begin{proof}[Proof of Theorem~\ref{ud:thm:unit}]
Widths zero, odd widths, and width two were excluded directly. At every
other width the first two interior words are forced by Section~\ref{ud:sec:forced}, and their
third-layer continuation is impossible by Section~\ref{ud:sec:interleaving}. Quarter-turn rotation
preserves the rule and finite support. No finite enumeration is a premise.
\end{proof}

\subsection{Growth consequence and attribution}
Jillhewar's results~\cite[Section~7, signed-residue theorem and the axis/even-diagonal growth-four theorems]{Jillhewar2026} give $g\in4\Z_{>0}$ and
$g\ge2\max(|d_x|,|d_y|)$. For $g=4$, both coordinates have absolute value at
most two. Since heading restoration makes $P$ even, $d_x+d_y$ is even.
Up to quarter-turns, the only nonzero drift types are $(1,-1)$, $(2,0)$ and
$(2,-2)$. Theorem~\ref{ud:thm:unit} excludes the first; the cited all-period
axis and even-diagonal growth-four theorems exclude the others. This recovers
$g\ge8$. The supplemental periodic-growth proof provides an attributed reconstruction
of those external dependencies. Neither these earlier algebraic results nor
the growth corollary is claimed as new here. The analytic unit-diagonal proof
uses only the preceding stable-rank and interleaving arguments and does not rely on those other exclusions.

\subsection{Independent computational checks}
The supplement's \texttt{reproduce.py} reruns a complete-row Python enumerator,
an independently written C++ vector-geometry enumerator, and witness audits.
They agree on $4/124/12\,756$ local candidates and $2/2/0$ survivors. A separate
implementation validates all $12\,880$ rejected branches, $32\,766$ short binary
words, $32$ small count domains, eight compass port cases, two corrupt-witness
negative controls, and $936$ genuine P104 prefix controls retaining the full
$(2,-2)$ quotient. An additional analytic audit verifies the displayed first
return paths, the forced words~\eqref{ud:eq:B}, and the interleaving implication
on every candidate in the finite enumeration. These checks corroborate the
local identities; the symbolic forcing and contradiction above establish the
all-period exclusion.
